\chapter{Contexto}

\begin{quotation} [Empresario] {Paul Allen (1953--2018)}
Me gusta crear nuevas ideas, trabajando en nuevos proyectos creativos.
\end{quotation}

\begin{abstract}
Comenzar un proyecto no es sencillo, lo primero que hemos de hacer es proveer de contexto para evaluar dónde nos encontramos.
\end{abstract}

\section{Hacia una Gestión Eficiente y Justa de Cambios de Grupo}

La implementación de un sistema de permutas surge de la necesidad de ofrecer una mayor flexibilidad a los estudiantes en la organización de sus horarios académicos. Con el actual proceso de asignación de grupos, los estudiantes a menudo se encuentran situaciones en las que su horario no se ajusta bien a sus responsabilidades personales, laborales u otras actividades académicas. Esto puede generar dificultades en el equilibrio entre el estudio y otros compromisos, afectando tanto el rendimiento académico como el bienestar personal de los estudiantes.

Actualmente, el proceso de permutas se realiza de manera manual y descentralizada, lo que implica que cada estudiante debe buscar por su cuenta a otra persona con quien intercambiar su grupo. Este sistema informal no garantiza que todos los estudiantes que necesitan un cambio logren realizarlo, y puede ser lento e ineficiente, especialmente en períodos de alta demanda al inicio de cada cuatrimestre.

La implementación de un sistema de permutas automatizado permitiría simplificar este proceso, facilitando que los estudiantes encuentren coincidencias de manera más rápida y justa. Un sistema de este tipo podría ofrecer:

\begin{itemize}

\item  \textbf{Eficiencia en el proceso de cambio de grupo}. Al identificar automáticamente posibles permutas, el sistema ahorraría tiempo tanto a estudiantes como a la administración.

\item  \textbf{Accesibilidad y equidad}. Todos los estudiantes tendrían las mismas oportunidades de realizar cambios, sin depender de contactos personales o de esfuerzos individuales.
    
\item \textbf{Gestión centralizada y soporte}. La administración podría supervisar y apoyar el proceso mediante un sistema de gestión de incidencias, asegurando que cualquier problema pueda ser resuelto rápidamente y ofreciendo un servicio de ayuda continuo al estudiantado.

\end{itemize}


\section{Estado del arte}

\textbf{Introducción}

En este apartado se detallan las investigaciones previas que se han llevado a cabo sobre la implementación de sistemas digitales de gestión de permutas.

\textbf{Investigaciones consultadas}\\
Algunas universidades en España han implementado sistemas digitales de gestión de permutas para facilitar el cambio de grupo entre estudiantes y asegurar un proceso más eficiente y equitativo. Estos sistemas generalmente operan a través de plataformas en línea donde los estudiantes pueden solicitar permutas de manera rápida y organizada.\\

Un ejemplo, es la Universidad de Las Palmas de Gran Canaria (ULPGC), donde el proceso de permuta de asignaturas es un sistema formalmente implementado para permitir a los estudiantes realizar cambios en sus matrículas. Este proceso está diseñado para gestionar situaciones en las que los estudiantes deseen cambiar asignaturas de su plan de estudios, ya sea por conflicto de horarios, necesidad de reorganizar su carga académica, o por otras razones justificadas.\\
Este sistema facilita la presentación de solicitudes de permuta a través de la Sede Electrónica de la ULPGC, permitiendo que el trámite se realice de manera digital, sin necesidad de que los estudiantes acudan físicamente a las oficinas administrativas. Los estudiantes deben especificar las asignaturas que desean permutar, incluyendo información como códigos y créditos, y cualquier ajuste en las tasas de matrícula se realiza según los créditos implicados en la permuta.\\
Además, la ULPGC regula el proceso de permuta en función de los periodos de matrícula de cada cuatrimestre. Esto asegura que los estudiantes pueden reorganizar sus asignaturas según sus necesidades en cada etapa del curso académico, dentro de los plazos oficiales. \footnote{Procedimiento para la permuta de asignaturas, ULPGC. (\href{https://www.ulpgc.es/sites/default/files/ArchivosULPGC/adm_eet/tramite_de_permutas.pdf}{en la siguiente referencia})}\\ 

Otro caso, sería el sistema de permutas implementado por la Universidad Autónoma de México (UAM) se lleva a cabo completamente de manera digital a través de su módulo de información escolar.\\
Este sistema permite a los estudiantes intercambiar grupos sin necesidad de realizar trámites presenciales, facilitando el proceso y reduciendo la carga administrativa. Para realizar una permuta, ambos estudiantes deben ingresar al sistema, registrar los datos del otro estudiante, y confirmar la disponibilidad de plazas en los grupos solicitados. Esta plataforma mejora la eficiencia y accesibilidad del proceso, permitiendo que los estudiantes gestionen sus horarios de forma ágil.\footnote{Para consultar más detalles del proceso de gestión de permutas, UAM. (\href{https://www.ulpgc.es/sites/default/files/ArchivosULPGC/adm_eet/tramite_de_permutas.pdf}{en la siguiente referencia})}\\

\textbf{Conclusiones}

Los sistemas de permutas en universidades españolas y a nivel internacional están evolucionando rápidamente hacia plataformas digitales más eficientes. La digitalización de estos procesos no solo mejora la eficiencia administrativa, sino que también facilita un acceso más equitativo para todos los estudiantes.\\
La implementación de sistemas avanzados, como algoritmos de optimización, y el establecimiento de canales de soporte, son tendencias claves que podrían servir de modelo para futuras implementaciones.                                                                                                                                                                                              